\documentclass[book.tex]{subfiles}

\begin{document}

\chapter{Convolutional Neural Networks}

\label{chap:convnets}
\noindent\rule{11cm}{0.4pt} \\
\textbf{Prerequisites:}  \autoref{chap:ml_basics}, \autoref{chap:grad_descent}  \\
\textbf{Difficulty Level:} ** \\
\noindent\rule{11cm}{0.4pt}
\vspace{5mm}

Convolutional networks are designed to process grid structured inputs such as images more efficiently than fully connected networks. The core idea is to make use of translational invariance in images (e.g. a dog in a picture should still remain a dog if shifted in the picture). Convolutional networks achieve this by using the mathematical operation of convolution which is natively translation invariant.

Each convolution in a convolutional network has a set of learned weights called a filter. Each filter is applied to the entire input grid by sliding the filter over the image. The use of these filters reduces the number of parameters when compared with a fully connected network. The output of a convolutional layer is a new output grid with a possibly different set of dimensions.

\begin{figure}
    \centering
    % \includegraphics{figures/Differentiable Models/conv_nets/REPLACE_karpathy_cnn.jpeg}
    \includegraphics[width=0.9\textwidth]{figures/Differentiable Models/conv_nets/convolutional network241.jpg}
    \caption{A convolutional network accepts a grid of input with width, height, and depth. Each convolutional layer transforms an input grid into an output grid. }
    \label{fig:my_label}
\end{figure}

\section{The Convolution Operation}
In its most general form, convolution is an operation on two functions of a real$-$valued argument. To motivate the definition of convolution, the following two functions are shown as examples. 

Suppose a user is interested in tracking the location of a spaceship with a laser sensor. The laser sensor provides a single output $x(t)$, the position of the spaceship at time $t$. Both $x$ and $t$ are real$-$valued, there can be a different reading from the laser sensor at any instant in time.

Now suppose that the laser sensor is somewhat noisy. To obtain a less noisy estimate of the spaceship's position, an average of several measurements can be taken. Of course, more recent measurements are more relevant, so this should be a weighted average that gives more weight to recent measurements. This is done with a weighting function $w(a)$, where $a$ is the age of a measurement. If this is applied such that there is a weighted average operation at every moment, there is a new function s providing a smoothed estimate of the position of the spaceship: 
\begin{align}
s(t) &= \int x(a)w(t-a)da
\end{align} 
This operation is called convolution. The convolution operation is typically denoted with an asterisk:
\begin{align}
s(t) &= (x * w)(t)
\end{align}
\indent In general, convolution is defined for any functions for which the above integral is defined, and may be used for other purposes besides taking weighted averages. In convolutional network terminology, the first argument (in this example, the function $x$) to the convolution is often referred to as the input and the second argument (in this example, the function $w$) as the kernel. The output is sometimes referred to as the feature map. Usually, for data on a computer, time will be discretized, and the sensor will provide data at regular intervals. In the example above, it might be more realistic to assume that the laser provides a measurement once per second. The time index t can then take on only integer values. If it is now assumed that $x$ and $w$ are defined only on integer $t$, the discrete convolution is then defined as:
\begin{align}
s(t) &= (x*w)(t) = \sum^\infty_{a=-\infty} x(a)w(t-a)
\end{align}
In machine learning applications, the input is usually a multidimensional array (tensor) of data and the kernel is usually a multidimensional array (tensor) of parameters that are adapted by the learning algorithm. Because each element of the input and kernel must be explicitly stored separately, it is usually assumed that these functions are zero everywhere but the finite set of points for which we store the values. This means that in practice, the infinite summation can be implemented as a summation over a finite number of array elements. Since convolutions are used over more than one axis at a time, a larger dimension kernel $K$ is generally needed. The convolution for the example of a two$-$ dimensional image $I$ is

\begin{align}
S(i,j) &=(I*K)(i,j) = \sum_m \sum_n I(m,n)K(i-m,j-n)
\end{align}
A convolution is commutative thus the following is equivalently true:
\begin{align}
S(i,j) &= (K*I)(i,j)=\sum_m \sum_n I(i-m,j-n)K(m,n)
\end{align}
Usually the latter formula is more straightforward to implement in a machine learning library, because there is less variation in the range of valid values of $m$ and $n$.

% \indent The commutative property of convolution arises because we have \textbf{flipped} the kernel relative to the input, in the sense that as m increases, the index into the input increases, but the index into the kernel decreases. Many neural network libraries implement a related function called the cross-correlation, which is the same as convolution but without flipping the kernel:
% \begin{equation}
%     S(i,j)=(I*K)(i,j)= \sum_m \sum_n I(i+m,j+n)K(m,n)
% \end{equation}
% In the context of machine learning, the learning algorithm will learn the appropriate values of the kernel in the appropriate place, so an algorithm based on convolution with kernel flipping will learn a kernel that is flipped relative to the kernel learned by an algorithm without the flipping. It is also rare for convolution to be used alone in machine learning; instead convolution is used simultaneously with other functions, and the combination of these functions does not commute regardless of whether the convolution operation flips its kernel or not. Discrete convolution can be viewed as multiplication by a matrix. However, the matrix has several entries constrained to be equal to other entries. For example, for univariate discrete convolution, each row of the matrix is constrained to be equal to the row above shifted by one element. This is known as a Toeplitz matrix. In two dimensions, a doubly block circulant matrix corresponds to convolution. In addition to these constraints that several elements be equal to each other, convolution usually corresponds to a very sparse matrix (a matrix whose entries are mostly equal to zero). Any neural network algorithm that works with matrix multiplication and does not depend on specific properties of the matrix structure should work with convolution, without requiring any further changes to the neural network.\\
It may be necessary to skip over some positions of the kernel in order to reduce the computational cost (at the expense of not extracting the features as finely). This can be thought of as downsampling the output of the full convolution function. If only every s pixels in each direction in the output are sampled, then a downsampled convolution function c is defined such that:
\begin{align}
Z_{i,j,k} &= c(K,V,s)_{i,j,k} = \sum_{l,m,n}[V_{l,(j-1)*s+m,(k-1)*s+n}K_{i,l,m,n}]
\end{align}
where $Z$ is the output, $K$ is the kernel, and $V$ is the observed data. $s$ in this equation is referred to as the stride of the downsampled convolution. It is possible to have separate strides for each direction of motion.

\subsection{Motivating the Definitions} 
Convolution leverages three important ideas that can help improve a machine learning system: sparse interactions, parameter sharing and equivariant representations. Moreover, convolution provides a means for working with inputs of variable size.  There is also a natural connection to differential equations.

%Traditional neural network layers use matrix multiplication by a matrix of parameters with a separate parameter describing the interaction between each input unit and each output unit. This means every output unit interacts with every input unit. 
Convolutional networks typically have sparse interactions (also referred to as sparse connectivity or sparse weights). This is accomplished by making the kernel smaller than the input. For example, when processing an image, the input image might have thousands or millions of pixels, however small, meaningful features such as edges can be detected with kernels that occupy only tens or hundreds of pixels. 

%This means that we need to store fewer parameters, which both reduces the memory requirements of the model and improves its statistical efficiency. It also means that computing the output requires fewer operations. In a deep convolutional network, units in the deeper layers may indirectly interact with a larger portion of the input. This allows the network to efficiently describe complicated interactions between many variables by constructing such interactions from simple building blocks that each describe only sparse interactions.\\
Parameter sharing refers to using the same parameter for more than one function in a model. In a fully connected neural net, each element of the weight matrix is used exactly once when computing the output of a layer. It is multiplied by one element of the input and then never revisited. As a synonym for parameter sharing, one can say that a network has tied weights, because the value of the weight applied to one input is tied to the value of a weight applied elsewhere. In a convolutional neural net, each member of the kernel is used at every position of the input (except perhaps some of the boundary pixels, depending on the design decisions regarding the boundary). The parameter sharing used by the convolution operation means that rather than learning a separate set of parameters for every location, only one set is learned. This further improves the memory requirements and statistical efficiency.

In the case of convolution, the particular form of parameter sharing causes the layer to have a property called equivariance to translation. 
%To say a function is equivariant means that if the input changes, the output changes in the same way. Specifically, a function f(x) is equivariant to a function g if f(g(x)) = g(f(x)). In the case of convolution, if g is any function that translates the input, i.e., shifts it, then the convolution function is equivariant to g. 
For example, let $I$ be a function giving image brightness at integer coordinates. Let g be a function mapping one image function to another image function, such that $I' = g(I)$ is the image function with 
\begin{align}I'(x, y) = I(x-1,y)\end{align}
This shifts every pixel of $I$ one unit to the right. If this transformation is applied to $I$, then convolution is applied, the result will be the same as if the convolution is applied to $I'$, then applied the transformation $g$ to the output. When processing time series data, this means that convolution produces a sort of timeline that shows when different features appear in the input. Similarly with images, convolution creates a 2-D map of where certain features appear in the input. Convolution is not naturally equivariant to some other transformations, such as changes in the scale or rotation of an image. Other mechanisms are necessary for handling these kinds of transformations.

\section{Formalizing the Convolutional Network}

Suppose that we have a $N \times N$ dimensional input $X_{ij}$. Suppose that we have a $m \times m$ weight filter $w_{ij}$. For simplicity, let us suppose that we take steps of size $1$. We can compute the $i,j$-th element of output $X'$ as follows

\begin{align}
    X'_{ij} &= \sum_{a=0}^{m-1} \sum_{b=0}^{m-1} w_{ab} X_{i+a,j+b}
\end{align}
Note the indices carefully. Our definition means that the size of the output $X'$ will be $(N - m +1)\times(N - m + 1)$ due to edge effects. Let's transform this update rule into Physika code
\begin{lstlisting}[language=python]
class conv2d(w: $\mathbb{R}[m, m]$):
  def $\lambda$(X: $\mathbb{R}[N, N]$) $\to$ $\mathbb{R}[N-m+1,N-m+1]$:
    X$'$ = zeros(N-m+1, n-m+1)
    for i j a b:
      X$'$[i, j] += w[a, b]*X[i, j]
    X$'$
\end{lstlisting}

Implementations of convolutions don't usually use nested for-loops in a high level language, since in practice, very efficient hardware operations are used under the hood to parallelize these calculations to the degree possible on modern hardware.

\subsection{Pooling Layers}
A typical layer of a convolutional network consists of three stages. In the first stage, the layer performs several convolutions in parallel to produce a set of linear activations. In the second stage, each linear activation is run through a nonlinear activation function, such as the rectified linear activation function. This stage is sometimes called the detector stage. In the third stage, we use a pooling function to modify the output of the layer further.

A pooling function replaces the output of the net at a certain location with a summary statistic of the nearby outputs. For example, the max pooling operation reports the maximum output within a rectangular neighborhood. Other popular pooling functions include the average of a rectangular neighborhood, the $L^2$ norm of a rectangular neighborhood, or a weighted average based on the distance from the central pixel. In all cases, pooling helps to make the representation become approximately invariant to small translations of the input. Invariance to translation means that if the input is translated by a small amount, the values of most of the pooled outputs do not change. This is important for detecting whether or not a feature is present and not caring about its exact location (i.e. identifying a human face). Pooling over spatial regions produces invariance to translation, but the pool is applied over the outputs of separately parameterized convolutions, the features can learn which transformations to become invariant to.

Since pooling summarizes the responses over a whole neighborhood, it is possible to use fewer pooling units than detector units, by reporting summary statistics for pooling regions spaced $k$ pixels apart rather than $1$ pixel apart. This improves the computational efficiency of the network because the next layer has roughly $k$ times fewer inputs to process. This reduction in the input size can also result in improved statistical efficiency and reduced memory requirements for storing the parameters.

For many tasks, pooling is essential for handling inputs of varying size. For example, if the goal is to classify images of variable size, the input to the classification layer must have a fixed size. This is usually accomplished by varying the size of an offset between pooling regions so that the classification layer always receives the same number of summary statistics regardless of the input size. For example, the final pooling layer of the network may be defined to output four sets of summary statistics, one for each quadrant of an image, regardless of the image size. 

\section{Deep Convolutional Networks}

One of the most powerful parts of a convolutional architecture is composability. For example, we can chain a series of convolutions together that performs the following series of shape changes by chaining together $k$ convolutional layers
\begin{align}
\mathbb{R}^{N \times N} \to \mathbb{R}^{(N-m+1)\times(N-m+1)}  \to \dotsc \to \mathbb{R}^{(N-km+1)\times(N-km+1)}
\end{align}
The shapes in the sequence above are shrinking continuously, placing a natural limit on the depth of any such network. Is there a way to prevent shrinkage? One trick is to perform zero-padding which pads the original array with zeros before convolving.
\begin{lstlisting}
def pad(X: $\mathbb{R}[N, N]$, m) $\to$ $\mathbb{R}[N+m-1,N+m-1]$:
  X' = zeros(N+m-1, N+m-1)
  X'[:N, :N] = X
  X'
\end{lstlisting}
By zero-padding before applying convolution, we get the following sequence of shape transformations
\begin{align}
\mathbb{R}^{N \times N} \to \mathbb{R}^{(N+m-1)\times(N+m-1)} \to
\mathbb{R}^{N \times N}
\end{align}
We can chain as many of these padded convolutions in a sequence as we would like. Convolutions can also be applied to 1 dimensional inputs, 3 dimensional inputs or to arbitrary $d$ dimensional input. We simply need to change the number of nested for-loops we perform.

\section{Residual Networks}

The residual network is a common motif in convolutional (and non-convolutional) architectures that we briefly introduce here. It provides a mechanism to preserve a copy of the original input to the layer in addition to applying a learned transformation.

\begin{marginfigure}
    \centering
    % \includegraphics{figures/Differentiable Models/conv_nets/residual.jpeg}
    \includegraphics{figures/Differentiable Models/conv_nets/residual_0.png}
    \caption{A residual layer preserves its input in addition to applying convolutions.}
    \label{fig:my_label}
\end{marginfigure}

\begin{lstlisting}[language=python]
class Residual(sublayer: Module, M: $\mathbb{N}$, dropout: $\mathbb{R}$):
    def $\lambda$(tensors: [Tensor]) $\to$ Tensor:
      tensors[0] + Dropout(sublayer(tensors))    
\end{lstlisting}
%    norm = LayerNorm(M)
%    dropout = Dropout(dropout)
\end{document}